% JETSPIN jet insertion documentation

\section{Jet insertion\label{subsec:Jet-insertion}}

The jet insertion at the nozzle is modeled as follows. For the sake
of simplicity, let us consider a simulation which starts with only
two bodies: a single mass-less point fixed at $x=0$ representing
the spinneret nozzle, and a bead modeling an initial jet segment of
mass $m_{i}$ and charge $e_{i}$ located at distance $l_{step}$
from the nozzle along the $x$ axis. Here, $l_{step}$ denotes the
length step used to discretize the jet in a sequence of beads. The
starting jet bead is assumed to have a cross-sectional radius $a_{0}$,
defined as the radius of the filament at the nozzle before the stretching
process. Applying the condition that the volume of the jet is conserved,
the relation $\pi a_{i}^{2}\,l_{i}=\pi a_{0}^{2}l_{step}$ is valid
for any $i-th$ bead. Furthermore, the starting jet bead has an initial
velocity $\upsilon_{s}$ along the $x$ axis equal to the bulk fluid
velocity in the syringe needle. Once this traveling jet bead is a
distance of $2\cdot l_{step}$ away from the nozzle, a new jet bead
(third body) is placed at distance $l_{step}$ from the nozzle along
the straight line joining the two previous bodies. Let us now label
$i-1$ the farthest bead from the nozzle, and $i$ the last inserted
bead. The $i-th$ bead is inserted with the initial velocity $\upsilon_{i}=\upsilon_{s}+\upsilon_{d}$,
where $\upsilon_{d}$ denotes the dragging velocity computed as

\begin{equation}
\upsilon_{d}=\frac{\upsilon_{i-i}-\upsilon_{s}}{2}.
\end{equation}

Here, the dragging velocity should be interpreted as an extra term
which accounts for the drag effect of the electrospun jet on the last
inserted segment. Note that the actual dragging velocity definition
was chosen in order to not alter the strain velocity term $\left(1/l_{i-1}\right)\cdot\left(dl_{i-1}/dt\right)$
of Eq \ref{eq:huang-garcia} before and after the bead insertion.
